\documentclass[10pt,a4paper]{amsart}
\usepackage[margin=19mm]{geometry}
\usepackage{amsmath,amssymb,mathtools}
\usepackage[scr=rsfs,cal=euler]{mathalfa}
\usepackage{xfrac}
\usepackage[unicode,hidelinks,bookmarks=false]{hyperref}
\newcommand{\ie}{i.e.\ }
\newcommand{\cf}{cf.\ }
\newcommand{\resp}{respectively\ }
\newcommand{\Subsection}{\subsection}
\providecommand{\nobreakdash}{\nobreak\hbox{-}}
\setlength{\emergencystretch}{2em}
\multlinegap0pt
\usepackage{mathtools}
\usepackage{amssymb}
\usepackage[shortlabels]{enumitem}
\usepackage{tikz}
\newtheorem{remark}{Remark}
\newtheorem{lemma}{Lemma}
\usepackage{cleveref}
\crefname{remark}{Remark}{Remarks}
\crefname{lemma}{Lemma}{Lemmas}
\newcommand{\R}{\mathbb{R}}
\title{Ergodicity of Langevin Dynamics and its Discretizations for Non-smooth Potentials with Linearly Growing Drift}
\author{Lorenz Frühwirth \and Andreas Habring}
\date{Source: arXiv:2411.12051v2 (CC BY 4.0)}
\begin{document}
\maketitle

\noindent\emph{Source and licence.} Ergodicity of Langevin Dynamics and its Discretizations for Non-smooth Potentials with Linearly Growing Drift. Version arXiv:2411.12051v2 (CC BY 4.0), distributed by arXiv under the Creative Commons Attribution 4.0 International licence, http://creativecommons.org/licenses/by/4.0/, as recorded in the arXiv licence metadata for this exact version and shown on the abstract page https://arxiv.org/abs/2411.12051v2. Full licence notice: \texttt{licenses/arxiv-2411.12051-CC-BY-4.0.txt}.\par\smallskip

\noindent\textit{Editorial scope.} Counted unit: the article's lemma lemma:lyapunov (source0.tex lines 632--640). The article's complete original proof of it is delivered with it (source0.tex lines 641--670, one \texttt{proof} environment of 30 lines). No mathematics is written, repaired or re-derived: the statement and the proof are the article's own lines, in the article's own notation, and no retained line is substituted or re-flowed.  Retained uncounted as the source-local context the proof needs: lines 368--380.  Nothing was omitted from any retained span, so the package carries no omission record.  No retained line contains a citation, so the wrapper carries no bibliography and no bibliography entry.  No retained line contains a figure environment, an image-inclusion command or drawing code, so no image is delivered and the package declares no asset dependency.  Editorial formatting only, in addition: an AMS \texttt{amsart} preamble that restates the article's own declarations and macros used by the retained lines, namely the environments \texttt{remark}, \texttt{lemma} on separate counters, together with the standard packages those lines need.  The retained environments print in this document as Remark 1, Lemma 1 in order of appearance, whereas the article prints the same items with the numbering of its own full document.  The complete native source of the article as distributed in the arXiv e-print for this exact version is bundled unchanged as \texttt{sources/168717/source0.tex}.
\begin{remark}\
\label{rem:Ass_A}
\begin{itemize}
    \item Strongly convex functions on $\R^d$ are coercive and continuous. Thus, $U(x)$ admits a minimizer which we can assume to be at $x=0$ without loss of generality. (In the converse case, sample from the potential $\hat{U}(x) = U(x+\hat{x})$ with $\hat{x}$ the minimizer and afterwards, add $\hat{x}$ to the obtained samples.) 
    \item Since $U$ is a convex function it follows immediately that $U \in H^1_{\text{loc}}( \R^d)$. Moreover, by Alexandrov's theorem, $U$ is (twice) differentiable in the classical sense on a set whose complement has Lebesgue measure zero. On that set the classical derivative of $U$ coincides with the weak derivative. On the remaining zero-set where this is not the case, one may select an arbitrary subgradient.
    \item The linear growth of the subgradient together with convexity immediately implies quadratic growth of $U$. Indeed, we have for $g\in\partial U(x)$ and some $c>0$
    \begin{equation}
        | U(x) | \leq |U(0)| + |\langle g,x\rangle |\leq |U(0)| + \|g\|\|x\|\leq c(1+\|x\|^2).
    \end{equation}
    % On that set the classical derivative of $U$ coincides with the weak derivative. Using a second order Taylor expansion, we can also infer that strong convexity together with at most quadratic growth ensures that $\nabla U$ grows at most linearly. Thus, quadratic growth of $U$ is equivalent to linear growth of its (sub)differential, which is a formulation which will be used frequently in this article. 
    \item By $m$-strong convexity of $U$, there exists a convex function $c(x)$ such that $U(x) = c(x) + \frac{m}{2}\|x\|^2$. Therefore, since convex functions in turn are lower bounded by linear functions, all moments of $\pi$ are finite, that is, for any $p\geq 0$, $\int_{\R^d} \|x\|^p e^{-U(x)} dx  < \infty.$
\end{itemize}
\end{remark}

\begin{lemma}[Lyapunov drift condition]\label{lemma:lyapunov}
    For $f(x)=\|x\|^2$, there exists $\eta>0$ independent of $\tau$ such that the discrete scheme satisfies for $\tau$ small enough the Lyapunov drift condition 
    \begin{equation*}         
        \begin{aligned}
            \Delta f(x) \coloneqq \mathbb{E} \left[ f(X_1^\tau) \middle| X_0^\tau = x \right] - f(x)
            \leq -m\tau f(x) + \eta\tau,\quad x\in\R^d
        \end{aligned}   
    \end{equation*}
\end{lemma}

\begin{proof}
    First note that by strong convexity of $U$ and since $U$ is minimal at $x=0$ (\Cref{rem:Ass_A}), $\langle\theta(x),x\rangle\geq m\|x\|^2$. Further, by linear growth of $\theta(x)$, there exist $K,c\geq 0$ such that for $\|x\|\geq K$, $\|\theta(x)\|^2\leq c\|x\|^2$. Let $C = \overline{B_{K}(0)}$ be the closed ball with center $0$ and radius $K$. Further, let $x \in C^c = \R^d \setminus C$. Then we have
    \begin{equation}\label{eq:lyapunov1}
        \begin{aligned}
            \|x-\tau\theta(x)\|^2 = \|x\|^2 - 2\tau\langle x,\theta(x)\rangle + \tau^2\|\theta(x)\|^2
            \leq & \|x\|^2 - 2m\tau\|x\|^2 + \tau^2c\|x\|^2\\
            \leq &\|x\|^2 \left( 1 - \tau \left(2m-\tau c \right)\right)\\
            \leq &\|x\|^2 \left( 1 - \tau m\right)
        \end{aligned}
    \end{equation}
    where the last inequality holds for $\tau<\frac{m}{c}$. Thus, for $x\in C^c$ we have
    \begin{equation*}         
        \begin{aligned}
            \mathbb{E} \left[ f(X_1^\tau) \middle| X_0^\tau = x \right]
            = \mathbb{E} \left[ \|x-\tau\theta(x) + \sqrt{2\tau} W_1\|^2 \middle| X_0^\tau = x \right]
            = &\|x-\tau\theta(x)\|^2 + 2\tau \\
            \leq &(1-m\tau)f(x) + 2\tau.
        \end{aligned}   
    \end{equation*}
    On the other hand, for $x\in C$ we find, since $C$ is bounded
    \begin{equation*}         
        \begin{aligned}
            \mathbb{E} \left[ f(X_1^\tau) \middle| X_0^\tau = x \right]
            = \mathbb{E} \bigg[\|x-\tau\theta(x) + \sqrt{2\tau} W_1\|^2 \bigg| &X_0^\tau = x \bigg]
            \leq  \|x\|^2 + c(K)\tau \\
            \leq & (1-m\tau)\|x\|^2 + m\tau \|x\|^2 + c(K)\tau \\
            \leq & (1-m\tau)f(x) + c \tau.
        \end{aligned}   
    \end{equation*}
\end{proof}

\end{document}
